\documentclass{article}%
\usepackage[T1]{fontenc}%
\usepackage[utf8]{inputenc}%
\usepackage{lmodern}%
\usepackage{textcomp}%
\usepackage{lastpage}%
\usepackage{geometry}%
\geometry{margin=2.2cm}%
%
\usepackage{amsmath}%
\usepackage{amssymb}%
\usepackage{booktabs}%
\usepackage{longtable}%
\usepackage{array}%
\usepackage{hyperref}%
\hypersetup{hidelinks}%
\title{ITA 2026 {-} Física (claude\_sonnet\_5\_5) x Poliedro\\[4pt]\large Comparação com o gabarito Poliedro Resolve\thanks{Gerado com gpt-resolve: \url{https://github.com/lgabs/gpt-resolve}}}%
\date{10/11/2026}%
%
\begin{document}%
\normalsize%
\maketitle%
\section*{Resumo}\textbf{Nota média: 9,5/10} (21 de 22 itens com a resolução correta; cada item vale o mesmo dentro da questão).\par\medskip%
\begin{longtable}{@{}c c p{0.78\textwidth}@{}}\toprule Questão & Nota & Resumo \\ \midrule
1 & 10,0 & A questão 1 está correta e concorda com o gabarito Poliedro nos três itens: 3/3. Em (a), as duas resoluções dão (v$^2$/g)(cos$\theta$ - $\mu$ sen$\theta$)/(sen$\theta$ + $\mu$ cos$\theta$). Em (b), v$^2$/(g tan$\theta$) é o mesmo que (v$^2$/g)cotg$\theta$. Em (c), as duas dão v$^2$/($\mu$g). Conferi numericamente que as formas são equivalentes e que o exemplo da resolução está certo. O modelo físico é o mesmo nos dois: roda interna na iminência de perder o contato e atrito no limite na roda externa. A resolução justifica melhor que esse é o caso crítico, pelo argumento do setor convexo de direções possíveis da força total. Só há pontos menores: a hipótese comum de ignorar torques, já que o enunciado não dá a altura do CM nem a bitola; a interpretação do item (c) como deslizamento lateral; e uma descrição do sentido do atrito no diagrama do gabarito que contradiz as próprias equações dele. \\
2 & 10,0 & Resolução e gabarito concordam nos dois itens. Item a: as normais são algebricamente idênticas, já que M(L/2-d)/(L-2d) = M/2; conferi por torques em relação a A e a B e numericamente. Item b: os dois dão f\_A+f\_B = (M+m+m\_c)g sin$\theta$, para cima no plano, e concluem que a estática só determina a soma. A resolução acrescenta a repartição f\_i = N\_i tan$\theta$ como hipótese declarada (mesmo $\mu$, escorregamento iminente), e ela está correta nessa hipótese. O único ponto relevante para levar a um professor é a indeterminação do item b: não se sabe qual resposta a banca aceitaria para o atrito 'em cada apoio'. \\
3 & 10,0 & A resolução e o gabarito estão certos e concordam nos dois itens. No item a, a resolução dá R\_r = -M{[}$\omega$$^2$(L-2a)/2 + g cos$\theta${]} num eixo que vai do pivô para o CM. Isso equivale ao P\_R = Mg cos$\theta$ + M$\omega$$^2$(L-2a)/2 do gabarito, que aponta para a extremidade curta. No item b, as duas dão $\theta$ = arccos{[}-$\omega$$^2$(L-2a)/(2g){]}. Verifiquei numericamente, integrando a força centrífuga distribuída, que as duas abordagens dão o mesmo resultado: a da resolução, pelo movimento do CM no referencial inercial, e a do gabarito, no referencial girante. A resolução ainda acrescenta, corretamente, a condição de existência $\omega$$^2$(L-2a) $\leq$ 2g e a interpretação $\theta$ > 90$^\circ$. Não há divergência que afete a nota. As dúvidas listadas são de baixa importância (hipótese a < L/2, convenção de sinal, condição de existência). \\
4 & 10,0 & A resolução está certa nos dois itens e concorda com o Poliedro. a) As duas dão y$_{2}$ = 2b ou y$_{2}$ = -4b. b) Os termos estrela--fragmento são algebricamente iguais aos do gabarito: estrela no foco (c,0) e y$_{2}$ = 2b. A resolução acrescenta a interação mútua -16Gm$^2$/(25$\surd$(a$^2$+49b$^2$)), que está correta, e mostra que ela é desprezível porque M ? m. Conferi tudo numericamente. O ponto em aberto é o enunciado: a hipótese de que o CM continua na órbita original não se sustenta com fragmentos tão afastados, e por isso o Poliedro sugere anular a questão. \\
5 & 5,0 & Item a: a resolução e o gabarito coincidem. W = (M1 - M2)$\cdot$g$\cdot$$\alpha$$\cdot$L$\cdot$70 = 700$\cdot$$\alpha$$\cdot$L$\cdot$(M1 - M2) J, com g = 10 m/s$^2$. Item b: as respostas diferem. A resolução deduz corretamente o calor fornecido com o termo de energia potencial (M1 + m/2)$\cdot$g$\cdot$$\alpha$$\cdot$L$\cdot$$\Delta$T, mas o despreza no resultado final, porque vale cerca de 4$\cdot$10$^{-}$$^{7}$ do resto para valores típicos. Com isso, entrega $\eta$ = (M1 - M2)g$\alpha$L / ($\rho$Vc + 3mR/M\_Cu). O gabarito do Poliedro mantém esses termos no denominador e dá a expressão exata pela 1ª lei. A resposta da resolução é uma aproximação de ordem dominante, numericamente indistinguível, mas formalmente incompleta. Como o enunciado é literal e sem dados numéricos, considero o gabarito certo no item b. A resolução provavelmente recebe nota parcial ou integral, conforme o critério do corretor. \\
6 & 10,0 & A resolução está correta nos três itens e coincide com o gabarito Poliedro: a) L = 85 cm; b) G = -16, imagem invertida; c) imagem final real, invertida e menor, a cerca de 5,2 cm depois da ocular, com A = -1/16 $\approx$ -0,063. Refiz as contas com frações exatas e confirmei todos os valores. A única diferença é de arredondamento: o gabarito dá -0,063 e o valor exato é -0,0625. Há também um pequeno deslize de redação na verificação final da resolução. Nenhum dos dois tem erro de física. \\
7 & 10,0 & A questão 7 está correta: nos dois itens a resolução bate com o gabarito Poliedro e as duas estão certas. a) I = I\_0 x\^{}\{N+1\} cos\^{}\{2N\}($\varphi$/N); b) r = sqrt(P/($\pi$ I\_0 x\^{}\{N+1\} cos\^{}\{2N\}($\varphi$/N))). Reli o enunciado no PDF original e ele confere com a transcrição. Conferi a recorrência numericamente. Só há imprecisões de redação nas verificações auxiliares da resolução (o limite N$\to$$\infty$), que não afetam a nota. \\
8 & 10,0 & A resolução está correta e coincide com o gabarito Poliedro nos dois itens. No item a, P\_min = 35 W: 36 kJ na bateria mais 6 kJ de aquecimento, em 1200 s. No item b, $\Delta$t\_max = 2,24$\pi$$^2$$\cdot$10\^{}{-}6 s, que é igual a (56/25)$\pi$$^2$$\cdot$10\^{}{-}6 s $\approx$ 2,2$\cdot$10\^{}{-}5 s. A resolução chega a esse valor por M = L = $\mu$0N$^2$A/? e o gabarito por N$\cdot$$\mu$0$\cdot$n$\cdot$$\pi$r$^2$. As duas expressões são idênticas porque n = N/?. Refiz as contas em Python e os valores conferem. Ficam só pontos menores: a leitura de $\Delta$t como tempo da rampa, adotada pelos dois, e a unidade de $\mu$0 impressa errada na prova. \\
9 & 10,0 & A resolução está correta nos dois itens e bate com o gabarito do Poliedro. No item a, tg$\theta$ = 1 dá $\theta$ = 45$^\circ$ e d = 600$\surd$2 nm $\approx$ 848,5 nm. No item b, sen$\theta$' = $\surd$2/3 dá x = 20$\surd$14/7 cm $\approx$ 10,7 cm. Conferi as contas em python e o enunciado contra o PDF original. A resolução ainda verifica que a aproximação de ângulos pequenos não vale e que, com d/$\lambda$ = $\surd$2, só existe o máximo de primeira ordem do laranja. Não encontrei erros. \\
10 & 10,0 & A resolução está correta e coincide com o gabarito Poliedro nos dois itens. a) 12,09375 anos ($\approx$ 12,1 anos), medidos na Terra. b) 9,675 anos ($\approx$ 9,7 anos), o tempo próprio da nave 2. O enunciado transcrito confere com a página 11 do PDF. Refiz as contas com frações exatas e elas batem (387/32 e 387/40). A resolução ainda traz uma checagem independente por transformação de Lorentz, que dá o mesmo valor. Não há divergência a levar ao professor. \\
\bottomrule\end{longtable}%
\section*{Principais pontos de dúvida}%
\begin{itemize}%
\item%
\textbf{Q2} \textit{(statement)} O item b pede o atrito 'em cada apoio', mas com dois apoios rígidos o problema é estaticamente indeterminado: só a soma f\_A+f\_B = (M+m+m\_c)g sin$\theta$ sai do equilíbrio. Não se sabe o que a banca do ITA espera: só a soma (como no Poliedro), uma repartição por hipótese (proporcional às normais, metades iguais) ou f = $\mu$N, com $\mu$ que não foi dado. A nota do item pode depender de como a banca trata essa indeterminação. A resolução apresenta em destaque (boxed) valores individuais que dependem de hipótese, embora deixe claro que a estática só dá a soma.%
\item%
\textbf{Q4} \textit{(statement)} A premissa central, de que o CM dos fragmentos continua na elipse original, só vale enquanto os fragmentos sentem praticamente o mesmo campo, isto é, sem efeitos de maré. Os dados dão |P$_{1}$P$_{2}$| = $\surd$(a$^2$+49b$^2$)/4, distância comparável às dimensões da órbita. Por isso o Poliedro considera o enunciado inconsistente e sugere anular a questão. Se a banca anular ou mudar o critério, a nota da questão muda. A resolução declara essa hipótese de forma explícita, mas não discute se ela é válida para essa separação tão grande.%
\item%
\textbf{Q5} \textit{(solution)} No item b, a resolução deduz Q\_H exato, com o termo (M1 + m/2)$\cdot$g$\cdot$$\alpha$$\cdot$L$\cdot$$\Delta$T, e depois o descarta. A justificativa é numérica, com valores típicos que o enunciado não fornece. A expressão final fica sem M1$\cdot$g$\cdot$$\alpha$$\cdot$L e sem m$\cdot$g$\cdot$$\alpha$$\cdot$L/2 no denominador. Um corretor do ITA pode exigir a expressão completa, como a do Poliedro. Vale perguntar ao professor se essa aproximação explícita receberia nota cheia ou parcial.%
\end{itemize}%
\section*{Questão 1}%
\href{https://poliedroresolve.sistemapoliedro.com.br/vestibulares/ita/2026/ita-2-fase-dia-3-30-10-2025-2-fase-dia-3/dissertativa-1-fisica-2-fase-dia-3-ita-2026}{Resolução do gabarito}\par\smallskip%
\begin{longtable}{@{}p{0.06\textwidth} p{0.33\textwidth} p{0.33\textwidth} p{0.1\textwidth} p{0.08\textwidth}@{}}\toprule Item & Resolução & Gabarito & Veredito & Certo \\ \midrule
a & R\_min = v$^2$(1 - $\mu$ tan$\theta$)/{[}g($\mu$ + tan$\theta$){]} = (v$^2$/g)(cos$\theta$ - $\mu$ sen$\theta$)/(sen$\theta$ + $\mu$ cos$\theta$) = v$^2$/{[}g tan($\theta$ + arctan $\mu$){]} & r\_min = (v$^2$/g)(cos$\theta$ - $\mu$ sen$\theta$)/(sen$\theta$ + $\mu$ cos$\theta$) & igual & ambos \\
b & R\_min = v$^2$/(g tan$\theta$) & r\_min = (v$^2$/g)$\cdot$cotg $\theta$ & equivalente & ambos \\
c & R\_min = v$^2$/($\mu$ g) & r\_min = v$^2$/($\mu$g) & igual & ambos \\
\bottomrule\end{longtable}%
\begin{itemize}%
\item%
\textbf{a)} A resolução escreve o resultado exatamente na forma do gabarito (segunda forma do boxed). As outras formas saem dividindo numerador e denominador por cos$\theta$ e são idênticas (conferido numericamente). O modelo físico é o mesmo: roda interna na iminência de perder o contato (N\_i = f\_i = 0), atrito da roda externa no limite $\mu$N\_E ao longo da geratriz, para baixo e para dentro, e equilíbrio das forças. A resolução ainda justifica, pelo argumento do setor convexo, por que esse é o caso crítico.%
\item%
\textbf{b)} v$^2$/(g tan$\theta$) = (v$^2$/g)cotg$\theta$. Os dois fazem $\mu$ = 0 em (a). A resolução também deduz direto: N cos$\theta$ = Mg e N sen$\theta$ = Mv$^2$/R.%
\item%
\textbf{c)} Os dois fazem $\theta$ = 0 em (a). A resolução também deduz direto: Mv$^2$/R $\leq$ $\mu$Mg.%
\end{itemize}%
\paragraph{Pontos de dúvida}%
\begin{itemize}%
\item%
\textit{(statement)} Nem a resolução nem o gabarito usam equilíbrio de torques. Eles leem "o eixo permanece na horizontal" como "não há tombamento", e o enunciado não dá a altura do centro de massa nem a bitola. Em um trem real, a perda de contato da roda interna também depende do momento da força centrífuga aplicada no CM. A leitura é razoável e é a mesma nos dois.%
\item%
\textit{(statement)} No item (c), com rodas cilíndricas e sem friso, nenhuma roda perde o contato. O que acontece é o deslizamento lateral do trem. Os dois tratam o caso como o limite $\theta$ = 0 de (a), e a resolução avisa que aqui a roda "perde o trilho" por deslizamento.%
\item%
\textit{(key)} Segundo a transcrição do diagrama do gabarito, Fat\_E aponta "para baixo/para fora". As equações (I) e (II) do próprio gabarito (+Fat\_E cos$\theta$ no sentido centrípeto e -Fat\_E sen$\theta$ na vertical) exigem atrito para baixo e para dentro, rumo ao centro da curva, que é o sentido coerente com a geometria do cone. Deve ser erro da descrição da imagem e não afeta o resultado. Há também o erro de digitação "$\mu$n\_Esen$\theta$" no lugar de "$\mu$N\_Esen$\theta$".%
\item%
\textit{(solution)} Problemas de redação na resolução. "Não há rolamento do conjunto" quer dizer que não há rotação de rolagem nem tombamento, mas pode ser confundido com o rolamento das rodas. Em "Para $\theta$ pequeno, R\_min(a) < R\_min(b), R\_min(c)", a restrição a $\theta$ pequeno é desnecessária: a desigualdade vale para quaisquer $\theta$, $\mu$ > 0 com $\theta$ + arctan $\mu$ < 90$^\circ$. O exemplo numérico confere (226 m, 1029 m e 300 m).%
\item%
\textit{(solution)} A leitura da figura (raio da roda maior do lado interno, de modo que as normais apontam para o centro da via) coincide com a figura original e com o diagrama do gabarito. Mesmo assim, o resultado depende dessa leitura: com o cone invertido, o sinal de $\theta$ mudaria.%
\end{itemize}%
\section*{Questão 2}%
\href{https://poliedroresolve.sistemapoliedro.com.br/vestibulares/ita/2026/ita-2-fase-dia-3-30-10-2025-2-fase-dia-3/dissertativa-2-fisica-2-fase-dia-3-ita-2026}{Resolução do gabarito}\par\smallskip%
\begin{longtable}{@{}p{0.06\textwidth} p{0.33\textwidth} p{0.33\textwidth} p{0.1\textwidth} p{0.08\textwidth}@{}}\toprule Item & Resolução & Gabarito & Veredito & Certo \\ \midrule
a & N\_A = g cos$\theta$ {[}M/2 + (m(L-d-a) + m\_c(b-d))/(L-2d){]}; N\_B = g cos$\theta$ {[}M/2 + (m(a-d) + m\_c(L-b-d))/(L-2d){]} & N\_A = {[}m(L-a-d) + M(L/2-d) + m\_C(b-d){]}$\cdot$g$\cdot$cos$\theta$/(L-2d); N\_B = {[}m(a-d) + M(L/2-d) + m\_C(L-b-d){]}$\cdot$g$\cdot$cos$\theta$/(L-2d) & equivalente & ambos \\
b & f\_A + f\_B = (M+m+m\_c) g sin$\theta$, para cima ao longo do plano. A resolução afirma que a estática só determina a soma (o problema é hiperestático). Sob a hipótese adicional de mesmo $\mu$ nos dois apoios e escorregamento iminente (é a repartição que exige o menor $\mu$), dá f\_A = N\_A tan$\theta$ e f\_B = N\_B tan$\theta$, com $\mu$ $\geq$ tan$\theta$. & Fat\_A + Fat\_B = (m+M+m\_C) g sen$\theta$, as duas paralelas ao plano e para cima. Não há dados para obter cada atrito separadamente: o problema é indeterminado nesse ponto. & equivalente & ambos \\
\bottomrule\end{longtable}%
\begin{itemize}%
\item%
\textbf{a)} As expressões são idênticas, porque M(L/2-d)/(L-2d) = M/2. Refiz o problema por torques em relação a A e a B e pela soma das forças perpendiculares ao plano, e conferi numericamente com valores aleatórios dos parâmetros: resolução e gabarito coincidem (diferença da ordem de 1e{-}15), e N\_A+N\_B = (M+m+m\_c)g cos$\theta$. Os braços também estão certos: N\_B e os pesos medidos a partir de A (L-2d, L/2-d, a-d, L-b-d); N\_A e os pesos medidos a partir de B (L-2d, L/2-d, L-d-a, b-d).%
\item%
\textbf{b)} A soma e o sentido coincidem, e a conclusão principal também: só a soma é determinada pela estática. Isso acontece porque os atritos e as componentes paralelas dos pesos estão sobre o eixo da barra e não geram torque. A repartição f\_i = N\_i tan$\theta$ que a resolução acrescenta vem declarada como hipótese e está correta dentro dela. De f\_i $\leq$ $\mu$N\_i, somando, sai $\mu$ $\geq$ tan$\theta$; com $\mu$ = tan$\theta$, as duas desigualdades viram igualdades. Conferi numericamente que (N\_A+N\_B)tan$\theta$ = (M+m+m\_c)g sin$\theta$. Não contradiz o gabarito: é um complemento condicional.%
\end{itemize}%
\paragraph{Pontos de dúvida}%
\begin{itemize}%
\item%
\textit{(statement)} O item b pede o atrito 'em cada apoio', mas com dois apoios rígidos o problema é estaticamente indeterminado: só a soma f\_A+f\_B = (M+m+m\_c)g sin$\theta$ sai do equilíbrio. Não se sabe o que a banca do ITA espera: só a soma (como no Poliedro), uma repartição por hipótese (proporcional às normais, metades iguais) ou f = $\mu$N, com $\mu$ que não foi dado. A nota do item pode depender de como a banca trata essa indeterminação. A resolução apresenta em destaque (boxed) valores individuais que dependem de hipótese, embora deixe claro que a estática só dá a soma.%
\item%
\textit{(solution)} A frase 'é necessário (e, com essa repartição, suficiente) que $\mu$\_e $\geq$ tan$\theta$ em cada apoio' só vale sob a hipótese de mesmo $\mu$ nos dois apoios. Com coeficientes diferentes, a condição necessária e suficiente é $\mu$\_A N\_A + $\mu$\_B N\_B $\geq$ (M+m+m\_c)g sin$\theta$, e um dos apoios pode ter $\mu$ < tan$\theta$.%
\item%
\textit{(solution)} Resolução e gabarito usam o modelo de barra fina: pesos, normais e atritos sobre a mesma reta. Na figura, os apoios são pequenos calços sob a barra e o bloco m fica sobre ela. Se essas alturas fossem consideradas, as componentes paralelas ao plano (atrito e mg sin$\theta$) gerariam torque e alterariam N\_A e N\_B. É a hipótese usual, e o enunciado não dá alturas.%
\item%
\textit{(key)} O gabarito não comenta a condição de validade N\_A $\geq$ 0 e N\_B $\geq$ 0 (a barra não pode se desprender de um apoio), que justifica o papel do contrapeso. A resolução menciona essa condição. Não afeta o resultado.%
\end{itemize}%
\section*{Questão 3}%
\href{https://poliedroresolve.sistemapoliedro.com.br/vestibulares/ita/2026/ita-2-fase-dia-3-30-10-2025-2-fase-dia-3/dissertativa-3-fisica-2-fase-dia-3-ita-2026}{Resolução do gabarito}\par\smallskip%
\begin{longtable}{@{}p{0.06\textwidth} p{0.33\textwidth} p{0.33\textwidth} p{0.1\textwidth} p{0.08\textwidth}@{}}\toprule Item & Resolução & Gabarito & Veredito & Certo \\ \midrule
a & R\_r = -M{[}$\omega$$^2$(L-2a)/2 + g cos$\theta${]}, com o eixo radial orientado do pivô para o CM. Módulo M{[}$\omega$$^2$(L-2a)/2 + g cos$\theta${]}, ao longo da haste, do CM para o pivô, ou seja, para o lado da extremidade curta. & P\_R = Mg cos$\theta$ + M$\omega$$^2$(L-2a)/2, ao longo da haste, apontando para a extremidade mais próxima do pivô (trecho de comprimento a). & equivalente & ambos \\
b & $\theta$ = arccos{[}-$\omega$$^2$(L-2a)/(2g){]}, com $\theta$ > 90$^\circ$, que só existe se $\omega$$^2$(L-2a) $\leq$ 2g. & $\theta$ = arccos{[}-$\omega$$^2$(L-2a)/(2g){]}, considerando a < L/2. & igual & ambos \\
\bottomrule\end{longtable}%
\begin{itemize}%
\item%
\textbf{a)} É a mesma força com convenções de sinal diferentes. A resolução escreve o valor com sinal num eixo que vai do pivô para o CM (por isso o sinal negativo). O gabarito dá o módulo e diz que a força aponta para o trecho curto. As duas formas também descrevem o mesmo sentido: do CM para o pivô, isto é, para a extremidade curta. Refiz a conta. O CM fica a d = L/2 - a do pivô. A aceleração centrípeta do CM é -$\omega$$^2$d r?, e a componente do peso ao longo de r? é Mg cos$\theta$. Então R\_r = -M($\omega$$^2$d + g cos$\theta$). Também integrei numericamente a força centrífuga distribuída ao longo da haste ($\int$$\lambda$$\omega$$^2$s ds, com s de -a a L-a), e os três valores coincidiram em casos aleatórios: a integral, a fórmula da resolução e a do gabarito. O gabarito concentra a massa de cada trecho no seu centro. Isso é exato aqui, porque a força centrífuga é linear em r. A força de Euler é perpendicular à haste e não entra na equação radial, como a resolução também observa ao dizer que $\alpha$ não aparece.%
\item%
\textbf{b)} As expressões são idênticas. Conferi numericamente: com M=1, L=1,2, a=0,2, $\omega$=2 e g=10, a fórmula dá $\theta$ $\approx$ 99,2$^\circ$, e nesse ângulo R\_r = 0 (até o erro de máquina). A resolução acrescenta a condição de existência $\omega$$^2$(L-2a) $\leq$ 2g e a interpretação de que $\theta$ > 90$^\circ$ (CM acima do nível do pivô). Essas observações estão corretas e não aparecem no gabarito.%
\end{itemize}%
\paragraph{Pontos de dúvida}%
\begin{itemize}%
\item%
\textit{(key)} O gabarito do item b não diz que o ângulo só existe se $\omega$$^2$(L-2a) $\leq$ 2g, isto é, se o argumento do arccos estiver em {[}-1, 0{]}. A resolução inclui essa condição. A falta dela no gabarito não muda a resposta, mas uma correção rigorosa poderia valorizar essa discussão.%
\item%
\textit{(statement)} O enunciado não diz que a < L/2. Isso só se tira da figura, e tanto o gabarito quanto a resolução assumem essa hipótese. Se fosse a > L/2, a forma com sinal -M{[}$\omega$$^2$(L/2-a) + g cos$\theta${]} continuaria valendo, desde que $\theta$ fosse medido no trecho de comprimento L-a. Nesse caso, porém, o sentido descrito em palavras ('para a extremidade curta') mudaria.%
\item%
\textit{(statement)} No item b, $\omega$ é tratado como um parâmetro independente de $\theta$, como o enunciado pede ('em função de M, L, a, $\omega$ e g'). No movimento real, $\omega$ e $\theta$ estão ligados pela conservação de energia, então nem todo par ($\omega$, $\theta$) é atingível numa mesma trajetória. Nenhum dos dois discute isso, e o enunciado não exige.%
\item%
\textit{(solution)} Convenção de sinal: a resposta principal da resolução é um valor com sinal (negativo) num eixo que vai do pivô para o CM, enquanto o gabarito dá o módulo e o sentido. A resolução explica a convenção e também dá o módulo e o sentido corretos. Na correção, isso deve ser aceito como equivalente.%
\item%
\textit{(key)} A resolução transcrita do gabarito tem trechos truncados ou com erro de digitação. Por exemplo, aparece 'P\_R + Mgcos$\theta$' onde deveria estar '=', e as linhas estão emendadas. Parece defeito da extração do texto, e a expressão final está correta.%
\end{itemize}%
\section*{Questão 4}%
\href{https://poliedroresolve.sistemapoliedro.com.br/vestibulares/ita/2026/ita-2-fase-dia-3-30-10-2025-2-fase-dia-3/dissertativa-4-fisica-2-fase-dia-3-ita-2026}{Resolução do gabarito}\par\smallskip%
\begin{longtable}{@{}p{0.06\textwidth} p{0.33\textwidth} p{0.33\textwidth} p{0.1\textwidth} p{0.08\textwidth}@{}}\toprule Item & Resolução & Gabarito & Veredito & Certo \\ \midrule
a & y$_{2}$ = 2b ou y$_{2}$ = -4b (CM em (4a/5, (b+y$_{2}$)/5) sobre a elipse original) & y$_{2}$ = 2b ou y$_{2}$ = -4b (com a mesma premissa do CM na órbita original; o Poliedro sugere anular a questão) & igual & ambos \\
b & U\_min = -(GMm/5){[}16/$\surd$(25a$^2$-24ac-15b$^2$) + 1/$\surd$(2a$^2$-2ac+3b$^2$){]} - 16Gm$^2$/(25$\surd$(a$^2$+49b$^2$)), com y$_{2}$ = 2b, estrela em (c,0) e c = $\surd$(a$^2$-b$^2$). O último termo é declarado desprezível (ordem m/M). & E\_pg,mín = -GMm$_{2}${[}4/$\surd$((3a/4-c)$^2$ + (b/4)$^2$) + 1/$\surd$((a-c)$^2$ + (2b)$^2$){]}, com y$_{2}$ = 2b, estrela em (c,0) e c = $\surd$(a$^2$-b$^2$) & equivalente & ambos \\
\bottomrule\end{longtable}%
\begin{itemize}%
\item%
\textbf{a)} Mesma premissa e mesma conta: x\_CM = 4a/5 e y\_CM = (b+y$_{2}$)/5. Com 16/25 + (b+y$_{2}$)$^2$/(25b$^2$) = 1, sai (b+y$_{2}$)$^2$ = 9b$^2$, ou seja, y$_{2}$ = 2b ou -4b. Conferi numericamente que os dois valores põem o CM sobre a elipse. A resolução também declara a hipótese de desprezar os efeitos de maré, que é justamente o ponto criticado pelo Poliedro.%
\item%
\textbf{b)} Os termos estrela--fragmento são idênticos. Com m$_{2}$ = m/5, tem{-}se (3a/4-c)$^2$ + b$^2$/16 = (25a$^2$-24ac-15b$^2$)/16, logo 4/d$_{1}$ = 16/$\surd$(25a$^2$-24ac-15b$^2$). Também (a-c)$^2$ + 4b$^2$ = 2a$^2$-2ac+3b$^2$. A única diferença é que a resolução inclui a interação mútua -Gm$_{1}$m$_{2}$/d, com d = $\surd$(a$^2$+49b$^2$)/4, e o gabarito não. Esse termo está correto e também é minimizado por y$_{2}$ = 2b (7b/4 < 17b/4). Em relação ao termo principal, ele é da ordem de m/M: nos testes numéricos ficou entre 10$^{-}$$^{5}$ e 10$^{-}$$^{4}$ quando m/M $\approx$ 10$^{-}$$^3$. Com M ? m, as duas respostas coincidem, e a própria resolução escreve a forma aproximada, que é exatamente a do gabarito. Conferi também, com valores aleatórios de a, b, m e M, que entre as quatro combinações (dois focos $\times$ dois valores de y$_{2}$) o mínimo é sempre o caso de foco (c,0) com y$_{2}$ = 2b, e que a expressão fechada da resolução reproduz o cálculo direto.%
\end{itemize}%
\paragraph{Pontos de dúvida}%
\begin{itemize}%
\item%
\textit{(statement)} A premissa central, de que o CM dos fragmentos continua na elipse original, só vale enquanto os fragmentos sentem praticamente o mesmo campo, isto é, sem efeitos de maré. Os dados dão |P$_{1}$P$_{2}$| = $\surd$(a$^2$+49b$^2$)/4, distância comparável às dimensões da órbita. Por isso o Poliedro considera o enunciado inconsistente e sugere anular a questão. Se a banca anular ou mudar o critério, a nota da questão muda. A resolução declara essa hipótese de forma explícita, mas não discute se ela é válida para essa separação tão grande.%
\item%
\textit{(key)} O gabarito calcula a 'energia potencial do sistema' só com os pares estrela--fragmento e omite a interação mútua -Gm$_{1}$m$_{2}$/d. A resolução inclui esse termo, o que é mais completo, e mostra que ele é desprezível (ordem m/M). As duas formas devem ser aceitas, mas vale confirmar com o professor qual se espera.%
\item%
\textit{(statement)} O enunciado não diz em qual foco está a estrela nem qual sinal de y$_{2}$ vale. A 'menor energia possível' foi lida como a minimização sobre essas duas ambiguidades (foco e sinal de y$_{2}$). A resolução e o gabarito fazem a mesma leitura e chegam à mesma escolha: foco (c,0) e y$_{2}$ = 2b.%
\end{itemize}%
\section*{Questão 5}%
\href{https://poliedroresolve.sistemapoliedro.com.br/vestibulares/ita/2026/ita-2-fase-dia-3-30-10-2025-2-fase-dia-3/dissertativa-5-fisica-2-fase-dia-3-ita-2026}{Resolução do gabarito}\par\smallskip%
\begin{longtable}{@{}p{0.06\textwidth} p{0.33\textwidth} p{0.33\textwidth} p{0.1\textwidth} p{0.08\textwidth}@{}}\toprule Item & Resolução & Gabarito & Veredito & Certo \\ \midrule
a & W = (M1 - M2)$\cdot$g$\cdot$$\alpha$$\cdot$L$\cdot$(90 - 20) = 700$\cdot$$\alpha$$\cdot$L$\cdot$(M1 - M2) (SI, g = 10 m/s$^2$) & W = 70$\cdot$(M1 - M2)$\cdot$g$\cdot$$\alpha$$\cdot$L & equivalente & ambos \\
b & $\eta$ = (M1 - M2)$\cdot$g$\cdot$$\alpha$$\cdot$L / ($\rho$\_H2O$\cdot$V$\cdot$c\_H2O + 3mR/M\_Cu) = 10$\cdot$$\alpha$$\cdot$L$\cdot$(M1 - M2) / ($\rho$\_H2O$\cdot$V$\cdot$c\_H2O + 3mR/M\_Cu) & $\eta$ = (M1 - M2)$\cdot$g$\cdot$$\alpha$$\cdot$L / {[}$\rho$\_H2O$\cdot$V$\cdot$c\_H2O + (m/M\_Cu)$\cdot$3R + M1$\cdot$g$\cdot$$\alpha$$\cdot$L + m$\cdot$g$\cdot$$\alpha$$\cdot$L/2{]} & diferente & gabarito \\
\bottomrule\end{longtable}%
\begin{itemize}%
\item%
\textbf{a)} É a mesma expressão. Com g = 10 m/s$^2$, 70$\cdot$g = 700. As duas resoluções somam o trabalho de erguer M1 por $\Delta$L = $\alpha$L$\Delta$T e descer M2 pela mesma distância. A resolução também mostra, corretamente, que o peso próprio da viga dá trabalho líquido nulo no ciclo.%
\item%
\textbf{b)} Pela 1ª lei, o calor fornecido no aquecimento é Q\_q = ($\rho$Vc + 3mR/M\_Cu)$\Delta$T + M1$\cdot$g$\cdot$$\Delta$L + m$\cdot$g$\cdot$$\Delta$L/2. A última parcela é a energia potencial ganha por M1 e pelo centro de massa da viga, que sobe $\Delta$L/2. A própria resolução escreve essa expressão exata, com o termo (M1 + m/2)g$\Delta$L, igual ao gabarito. Depois ela descarta esse termo por ser desprezível e entrega uma aproximação de ordem dominante. Conferi em Python com os valores típicos da resolução (M1 = 100 kg, M2 = 50 kg, m = 7 kg, L = 1 m, V = 10 L): o termo descartado vale cerca de 3,9$\cdot$10$^{-}$$^{7}$ do denominador, e $\eta$ muda só na 7ª casa significativa (1,905423$\cdot$10$^{-}$$^{7}$ contra 1,905422$\cdot$10$^{-}$$^{7}$). A física da resolução está certa e a aproximação foi declarada. Mas o enunciado é literal e não dá valores numéricos, então nada autoriza cortar termos. A resposta exata pedida é a do gabarito, e a da resolução é só uma aproximação dela.%
\end{itemize}%
\paragraph{Pontos de dúvida}%
\begin{itemize}%
\item%
\textit{(solution)} No item b, a resolução deduz Q\_H exato, com o termo (M1 + m/2)$\cdot$g$\cdot$$\alpha$$\cdot$L$\cdot$$\Delta$T, e depois o descarta. A justificativa é numérica, com valores típicos que o enunciado não fornece. A expressão final fica sem M1$\cdot$g$\cdot$$\alpha$$\cdot$L e sem m$\cdot$g$\cdot$$\alpha$$\cdot$L/2 no denominador. Um corretor do ITA pode exigir a expressão completa, como a do Poliedro. Vale perguntar ao professor se essa aproximação explícita receberia nota cheia ou parcial.%
\item%
\textit{(solution)} A resolução diz que o termo mecânico é menor que correções que o modelo já ignora, como c\_P - c\_V do cobre (da ordem de 0,7 J/(mol$\cdot$K), cerca de 80 J/K para cerca de 110 mol, contra cerca de 0,02 J/K do termo mecânico). Fisicamente faz sentido, mas mistura o modelo idealizado do enunciado (c\_V = 3R, viga incompressível) com correções reais. Dentro do modelo, a 1ª lei inclui o termo mecânico de forma exata.%
\item%
\textit{(key)} O gabarito inclui a energia potencial do próprio peso da viga (m$\cdot$g$\cdot$$\Delta$L/2) no calor fornecido. Isso está correto: o centro de massa sobe $\Delta$L/2 no aquecimento e desce no resfriamento. Mas é uma escolha sutil, e outras resoluções podem omitir esse termo, ou também o de M1. Não se sabe se o ITA cobra esses termos.%
\item%
\textit{(statement)} Hipóteses que as duas resoluções usam sem declarar no enunciado: capacidade térmica do recipiente desprezível; dilatação da água e trabalho contra a atmosfera desprezados; viga em equilíbrio térmico com a água; uso de c\_V = 3R mesmo com a viga dilatando sob pressão constante. A resolução também trata M\_Cu em kg/mol. São hipóteses razoáveis e não mudam a forma das respostas.%
\end{itemize}%
\section*{Questão 6}%
\href{https://poliedroresolve.sistemapoliedro.com.br/vestibulares/ita/2026/ita-2-fase-dia-3-30-10-2025-2-fase-dia-3/dissertativa-6-fisica-2-fase-dia-3-ita-2026}{Resolução do gabarito}\par\smallskip%
\begin{longtable}{@{}p{0.06\textwidth} p{0.33\textwidth} p{0.33\textwidth} p{0.1\textwidth} p{0.08\textwidth}@{}}\toprule Item & Resolução & Gabarito & Veredito & Certo \\ \midrule
a & L = f1 + f2 = 85 cm (luneta afocal, os focos da objetiva e da ocular coincidem) & L = f1 + f2 = 80 + 5 = 85 cm & igual & ambos \\
b & G = -f1/f2 = -16 (|G| = 16); imagem invertida & G = f1/f2 = 16 em módulo; com sinal, G = -16; imagem invertida & igual & ambos \\
c & p1' = -80 cm (imagem intermediária virtual); p2 = 165 cm; p2' = 165/32 $\approx$ 5,2 cm depois da ocular ($\approx$ 90 cm da objetiva); imagem final real, invertida, A = -1/16 $\approx$ -0,063, menor que o objeto & p1' = -80 cm; p2 = 165 cm; p2' $\approx$ 5,2 cm; A $\approx$ -0,063; imagem real, invertida, menor, a 5,2 cm da ocular & igual & ambos \\
\bottomrule\end{longtable}%
\begin{itemize}%
\item%
\textbf{a)} Mesmo valor e mesma justificativa: imagem final no infinito exige p2 = f2, e a imagem da objetiva fica em f1. Então L = 80 + 5,0 = 85 cm.%
\item%
\textbf{b)} Os dois dão o módulo 16, o sinal negativo e a orientação invertida. A dedução da resolução ($\theta$' = h/f2 com h = -f1$\cdot$$\theta$) é o caso D $\to$ $\infty$ da expressão do gabarito, G = (f1/f2)(1 + f2/D).%
\item%
\textbf{c)} Refiz a conta com frações exatas: p1' = -80 cm, p2 = 85 + 80 = 165 cm, p2' = 165$\cdot$5/160 = 165/32 = 5,156 cm, A1 = +2, A2 = -1/32, A = -1/16 = -0,0625. O gabarito chega a -0,063 porque arredonda p2' para 5,2 antes de multiplicar (2$\cdot$5,2/165 = 0,0630). É só arredondamento. A posição (5,2 cm após a ocular), o tipo (real), a orientação (invertida) e o tamanho (menor) coincidem.%
\end{itemize}%
\paragraph{Pontos de dúvida}%
\begin{itemize}%
\item%
\textit{(statement)} No item (c), tanto a resolução quanto o gabarito supõem que o tubo continua com L = 85 cm, sem novo ajuste de foco para o objeto próximo. O enunciado não diz isso, mas é a leitura natural, porque L foi pedido no item (a).%
\item%
\textit{(statement)} A frase 'imagens formadas a longas distâncias do observador' foi lida como imagem final no infinito (luneta afocal, D $\to$ $\infty$), que leva a L = f1 + f2 e |G| = f1/f2. Se a imagem final ficasse no ponto próximo (D = 25 cm), os valores seriam outros (G = 19,2 e L $\approx$ 84,2 cm). O enunciado favorece claramente o infinito, e o gabarito usa a mesma leitura.%
\item%
\textit{(solution)} Na verificação (ii), a resolução chama a imagem intermediária de 'objeto virtual{-}intermediário'. Antes ela tinha dito, corretamente, que essa imagem é virtual para a objetiva, mas funciona como objeto REAL para a ocular. É só um deslize de redação e não afeta as contas nem a conclusão.%
\item%
\textit{(key)} O gabarito arredonda p2' para 5,2 cm antes de calcular A e obtém -0,063. O valor exato é -1/16 = -0,0625. A diferença é desprezível.%
\end{itemize}%
\section*{Questão 7}%
\href{https://poliedroresolve.sistemapoliedro.com.br/vestibulares/ita/2026/ita-2-fase-dia-3-30-10-2025-2-fase-dia-3/dissertativa-7-fisica-2-fase-dia-3-ita-2026}{Resolução do gabarito}\par\smallskip%
\begin{longtable}{@{}p{0.06\textwidth} p{0.33\textwidth} p{0.33\textwidth} p{0.1\textwidth} p{0.08\textwidth}@{}}\toprule Item & Resolução & Gabarito & Veredito & Certo \\ \midrule
a & I = I\_0 x\^{}\{N+1\} cos\^{}\{2N\}($\varphi$/N) & I\_final = I0$\cdot$x\^{}(N+1)$\cdot$cos\^{}(2N)(?/N) & igual & ambos \\
b & R = sqrt( P / ($\pi$ I\_0 x\^{}\{N+1\} cos\^{}\{2N\}($\varphi$/N)) ) & r = sqrt( P / ($\pi$$\cdot$I0$\cdot$x\^{}(N+1)$\cdot$cos\^{}(2N)(?/N)) ) & igual & ambos \\
\bottomrule\end{longtable}%
\begin{itemize}%
\item%
\textbf{a)} As duas usam o mesmo modelo: I\_1 = x I\_0 no primeiro polarizador, porque a luz chega paralela ao eixo e o caso ideal seria I\_0. Depois, cada um dos N polarizadores seguintes multiplica por x cos$^2$$\beta$ (lei de Malus com perda x). Daí vem I = I\_0 x\^{}\{N+1\} cos\^{}\{2N\}($\varphi$/N). Conferi numericamente, iterando a recorrência, e o valor coincide com a fórmula fechada.%
\item%
\textbf{b)} Como o feixe é cilíndrico e homogêneo, P = I$\cdot$$\pi$r$^2$ e portanto r = sqrt(P/($\pi$I)), com o I do item a. A expressão é a mesma nas duas. A resolução ainda dá a forma equivalente (1/(x\^{}\{(N+1)/2\} cos\^{}N($\varphi$/N)))$\cdot$sqrt(P/($\pi$I\_0)), que está correta.%
\end{itemize}%
\paragraph{Pontos de dúvida}%
\begin{itemize}%
\item%
\textit{(solution)} Na verificação do limite N$\to$$\infty$, a resolução diz que I $\to$ I\_0 x\^{}\{N+1\}. Esse limite não é bem definido, porque x\^{}\{N+1\} $\to$ 0 quando x<1 e N cresce. O que tende a 1 é só o fator de Malus cos\^{}\{2N\}($\varphi$/N). A redação é imprecisa, mas fica numa checagem auxiliar e não afeta as respostas.%
\item%
\textit{(solution)} A hipótese cos $\beta$ > 0 ($\beta$ < 90$^\circ$) é desnecessária, porque cos\^{}\{2N\} é par. Além disso, a desigualdade cos $\beta$ < 1 só vale para $\beta$ $\neq$ 0, e o próprio texto trata $\varphi$ = 0 como caso particular. São detalhes de redação.%
\item%
\textit{(statement)} O enunciado deixa implícitas duas coisas: que o 'analisador' apenas mede a potência, sem atuar como mais um polarizador, e que o fator x também vale para o primeiro polarizador. A resolução e o gabarito adotam a mesma leitura (I\_1 = x I\_0, sem polarizador extra), que é a mais natural.%
\end{itemize}%
\section*{Questão 8}%
\href{https://poliedroresolve.sistemapoliedro.com.br/vestibulares/ita/2026/ita-2-fase-dia-3-30-10-2025-2-fase-dia-3/dissertativa-8-fisica-2-fase-dia-3-ita-2026}{Resolução do gabarito}\par\smallskip%
\begin{longtable}{@{}p{0.06\textwidth} p{0.33\textwidth} p{0.33\textwidth} p{0.1\textwidth} p{0.08\textwidth}@{}}\toprule Item & Resolução & Gabarito & Veredito & Certo \\ \midrule
a & P\_min = 35 W (E = 3,6$\cdot$10\^{}4 J na bateria + 6,0$\cdot$10\^{}3 J de calor = 4,2$\cdot$10\^{}4 J em 1200 s) & P\_min = 35 W (36000 J + 6000 J = 42000 J em 1200 s) & igual & ambos \\
b & $\Delta$t\_max = 14L/U = 2,24$\pi$$^2$$\cdot$10\^{}{-}6 s $\approx$ 2,2$\cdot$10\^{}{-}5 s (L = $\mu$0N$^2$A/? = 8$\pi$$^2$$\cdot$10\^{}{-}7 H, M = L, $\Delta$I = 14 A) & $\Delta$t\_max = (56/25)$\cdot$$\pi$$^2$$\cdot$10\^{}{-}6 s $\approx$ 2,2$\cdot$10\^{}{-}5 s & igual & ambos \\
\bottomrule\end{longtable}%
\begin{itemize}%
\item%
\textbf{a)} Recalculei: q = 2 Ah = 7200 C, E\_bat = 7200$\cdot$5 = 36000 J; Q = 100 g$\cdot$2,0 J/(g$^\circ$C)$\cdot$30 $^\circ$C = 6000 J; P = 42000/1200 = 35 W. A resolução usa m em gramas e o gabarito usa 0,1 kg com c = 2000 J/(kg$^\circ$C). É a mesma conta. A justificativa do 'mínimo' (perdas adicionais só aumentam a energia fornecida) está correta.%
\item%
\textbf{b)} 56/25 = 2,24, então os valores coincidem. O gabarito usa $\varepsilon$ = N$\cdot$$\mu$0$\cdot$n$\cdot$$\pi$r$^2$$\cdot$$\Delta$I/$\Delta$t. A resolução usa $\varepsilon$ = M$\cdot$$\Delta$I/$\Delta$t com M = L = $\mu$0N$^2$A/?. Como N/? = n, as duas expressões são a mesma. Conferido em Python: N = 40, n = 2000 m\^{}{-}1, r = 5$\cdot$10\^{}{-}3 m e L = 8$\pi$$^2$$\cdot$10\^{}{-}7 $\approx$ 7,90$\cdot$10\^{}{-}6 H dão $\Delta$t = 2,2108$\cdot$10\^{}{-}5 s pelas duas fórmulas. Para justificar o 'máximo', a resolução cita $\Delta$I $\leq$ 14 A, k $\leq$ 1 e o fato de a idealização superestimar L. O gabarito diz que um $\Delta$t maior induziria menos de 5 V. Os dois argumentos são compatíveis.%
\end{itemize}%
\paragraph{Pontos de dúvida}%
\begin{itemize}%
\item%
\textit{(statement)} Leitura de $\Delta$t: a prova diz que a corrente 'varia linearmente entre 0 e 14 A em um período de tempo $\Delta$t'. A resolução e o gabarito tomam $\Delta$t como o tempo da rampa de 0 a 14 A, com taxa 14/$\Delta$t. Se alguém entender $\Delta$t como o período de uma onda triangular completa (0$\to$14$\to$0), a taxa passa a ser 28/$\Delta$t e o resultado dobra ($\approx$4,4$\cdot$10\^{}{-}5 s). A redação favorece a leitura adotada, que é a mesma do gabarito, mas vale registrar.%
\item%
\textit{(statement)} Unidade de $\mu$0 impressa errada: as instruções da prova trazem 'N.A\^{}{-}1' em vez de N$\cdot$A\^{}{-}2, e o gabarito escreve 'N/A'. Isso não afeta o valor numérico.%
\item%
\textit{(solution)} A resolução acrescenta uma checagem (fem própria do primário de 5 V $\times$ corrente média de 7 A = 35 W, igual ao item a) e cita a correção de Nagaoka ($\approx$20\%). Isso vai além do que se pede e não altera o resultado. A checagem pressupõe bobinas ideais sem perdas e só a fase de subida da rampa. Serve como coerência, não como demonstração.%
\end{itemize}%
\section*{Questão 9}%
\href{https://poliedroresolve.sistemapoliedro.com.br/vestibulares/ita/2026/ita-2-fase-dia-3-30-10-2025-2-fase-dia-3/dissertativa-9-fisica-2-fase-dia-3-ita-2026}{Resolução do gabarito}\par\smallskip%
\begin{longtable}{@{}p{0.06\textwidth} p{0.33\textwidth} p{0.33\textwidth} p{0.1\textwidth} p{0.08\textwidth}@{}}\toprule Item & Resolução & Gabarito & Veredito & Certo \\ \midrule
a & d = 600$\surd$2 nm $\approx$ 8,5$\cdot$10$^{-}$$^{7}$ m ($\theta$\_laranja = 45$^\circ$) & d = 600$\surd$2 nm ($\approx$ 848,5 nm), com tg$\theta$ = 1 $\Rightarrow$ $\theta$ = 45$^\circ$ & igual & ambos \\
b & x\_V = 20$\surd$14/7 cm $\approx$ 10,7 cm & y' = 20$\surd$14/7 cm ($\approx$ 10,7 cm) & igual & ambos \\
\bottomrule\end{longtable}%
\begin{itemize}%
\item%
\textbf{a)} Os dois usam o mesmo caminho: tg$\theta$ = 20/20 = 1, então $\theta$ = 45$^\circ$, e d$\cdot$sen45$^\circ$ = 600 nm, o que dá d = 600$\surd$2 nm $\approx$ 848,5 nm. Refiz a conta em python e confere. O enunciado do PDF (p. 10) bate com a transcrição.%
\item%
\textbf{b)} sen$\theta$' = 400/(600$\surd$2) = $\surd$2/3 $\approx$ 0,4714 e $\theta$' $\approx$ 28,1$^\circ$. A resolução calcula tg$\theta$' = $\surd$(2/7) e depois x = 20$\cdot$tg$\theta$'. O gabarito resolve y'/$\surd$(D$^2$+y'$^2$) = $\surd$2/3 e chega a 7y'$^2$ = 2D$^2$. As duas formas são equivalentes e dão 20$\surd$14/7 $\approx$ 10,69 cm, valor conferido numericamente. A resolução também mostra, corretamente, que a aproximação de ângulos pequenos (13,3 cm) não vale aqui.%
\end{itemize}%
\paragraph{Pontos de dúvida}%
\begin{itemize}%
\item%
\textit{(statement)} Hipótese implícita, a mesma na resolução e no gabarito: o "ponto central" é o pé da perpendicular da grade ao solo, onde cai o máximo de ordem zero. Também se supõe a luz solar como feixe paralelo, ignorando o diâmetro angular do Sol, que é de cerca de 0,5$^\circ$. Isso não altera o resultado.%
\item%
\textit{(solution)} Problema só de redação: a resolução diz "Como H é de poucos centímetros" para justificar que o ângulo não é pequeno. O que importa é a razão x/H $\approx$ 1, e não o valor absoluto de H. O raciocínio continua correto.%
\end{itemize}%
\section*{Questão 10}%
\href{https://poliedroresolve.sistemapoliedro.com.br/vestibulares/ita/2026/ita-2-fase-dia-3-30-10-2025-2-fase-dia-3/dissertativa-10-fisica-2-fase-dia-3-ita-2026}{Resolução do gabarito}\par\smallskip%
\begin{longtable}{@{}p{0.06\textwidth} p{0.33\textwidth} p{0.33\textwidth} p{0.1\textwidth} p{0.08\textwidth}@{}}\toprule Item & Resolução & Gabarito & Veredito & Certo \\ \midrule
a & $\Delta$t\_Terra = 387/32 anos = 12,09375 anos $\approx$ 12,1 anos & t\_Terra = 10,75 + 1,34375 = 12,09375 anos $\approx$ 12,1 anos & igual & ambos \\
b & $\Delta$$\tau$$_{2}$ = 9,675 anos $\approx$ 9,7 anos (tempo próprio, $\Delta$t/$\gamma$$_{2}$ com $\gamma$$_{2}$ = 1,25) & t\_nave2 = 12,09375/(5/4) = 9,675 anos & igual & ambos \\
\bottomrule\end{longtable}%
\begin{itemize}%
\item%
\textbf{a)} Refiz a conta no referencial da Terra. A nave 1 chega a Sirius em t1 = 8,6/0,8 = 10,75 anos. Nesse instante a nave 2 está em 0,6$\cdot$10,75 = 6,45 al, a 2,15 al da nave 1. Sinal e nave se aproximam a 1,6c, então o encontro leva 2,15/1,6 = 1,34375 ano, e o total é 12,09375 anos (= 387/32, conferido com frações exatas em Python). A resolução usa a equação horária e confere com o mesmo raciocínio de velocidade relativa do gabarito.%
\item%
\textbf{b)} Os dois eventos (passagem pela fronteira e recepção do sinal) ocorrem sobre a nave 2. Por isso o intervalo medido por ela é o tempo próprio: 12,09375$\cdot$0,8 = 9,675 anos (= 387/40). A resolução também confere com uma transformação de Lorentz completa. A emissão fica em t' = 6,9875 anos e x' = 2,6875 al, e somando o tempo de trânsito da luz dá $\tau$ = 9,675 anos. Conferi esses valores numericamente.%
\end{itemize}%
\paragraph{Pontos de dúvida}%
\begin{itemize}%
\item%
\textit{(statement)} Por "tempo transcorrido na Terra" entende{-}se a coordenada temporal no referencial Terra/Sol, que é a leitura usada pelos dois. O tempo de trânsito da luz até a Terra e a distância entre a Terra e a fronteira do sistema Solar foram desprezados. Isso é implícito e razoável, já que o enunciado dá t0 "para um observador na Terra".%
\item%
\textit{(solution)} A resolução supõe que as duas naves passam pela fronteira no mesmo ponto (trajetórias paralelas e próximas) e desconsidera a separação transversal. Isso não altera os resultados, e o gabarito faz a mesma hipótese ao dizer que a nave 2 "está junto da nave 1 na borda do sistema solar".%
\end{itemize}%
\end{document}